% s-180-p180-m002; original PDF 180
\Aphorism{s-180-p180-m002}{269}
Wille und Welle, Geist und Giessen, Seele und See: der Deutsche deutet sich den inneren Menschen hydraulisch, nicht pneumatisch. Bei ihm kommen auch die neugeborenen Kinder aus dem Teich und werden durch einen Wasservogel, den Storch, ins Haus gebracht; Quellen, Brunnen, Weiher sind des Deutschen geweihte Stätten und Märchen-Schauplätze. Die Mystik des Windes, von der das Johannesevangelium spricht, die Spiele und Zauber von Luft und Licht sagen dem Deutschen weniger — dazu war sein Klima nie trocken genug.
\par

Auch in der deutschen Ernährung spielt das Flüssige die Rolle, die es im deutschen Wetter spielt: der Deutsche ist kein \Foreign{animal animans}, sondern \Foreign{bibulum}, kein athmendes, sondern ein trinkendes Thier, das sich selbst die atmosphärische Luft verflüssigt. Lange dursten können, wie der Löwe in der Wüste, ist keine deutsche Tugend; lange dursten müssen ist deutsche „Noth“. Bei den Deutschen und ihren Stammesverwandten blüht auch jede Art Wassersport und Wasserheilkunde, bis zum offenbaren Unsinn; in England wird selbst das Temperament zu Wasser. Die romanischen Völker leben dagegen in einer rathsamen Unsauberkeit und fahren nicht gleich mit dem Scheuerlappen dazwischen, wenn sich in gesunder Luft eine gesunde Patina ansetzen will.
\par

